\begin{abstract}
Text-based emotion detection from social media has grown in importance as platforms such as Twitter generate large volumes of emotionally expressive informal text. This study proposes MindTrace, a comparative evaluation of four classification models --- SVM, XGBoost, CNN, and BiLSTM --- applied to a 416,809-sample Twitter corpus annotated across six emotion categories: anger, joy, sadness, fear, surprise, and love. Existing comparative studies vary widely in dataset selection, preprocessing methods, and evaluation metrics, making direct performance comparisons unreliable. This study addresses that gap by fixing all experimental conditions across all four model families on the same corpus. A standardised NLP preprocessing pipeline handles tokenisation, stopword removal with negation preservation, lemmatisation, and class balancing. ML models use TF-IDF vectorisation; DL models use trainable embedding layers. Performance is assessed using macro-F1, weighted-F1, and class-wise confusion matrices to account for significant class imbalance (imbalance ratio $\approx 9.5{:}1$). The findings are expected to provide empirically grounded guidance for architecture selection in emotion classification tasks, with applications in social media analytics, mental health monitoring, and emotion-aware AI systems.
\end{abstract}

\noindent\textbf{Keywords:} emotion classification, natural language processing, machine learning, deep learning, BiLSTM, CNN, SVM, XGBoost, TF-IDF, Twitter, multi-class classification, social media analytics
